\documentclass[conference]{IEEEtran}
% The reference PDF embeds Times New Roman / Times faces.
% Under pdflatex, mathptmx keeps this report in the same family.
\usepackage[T1]{fontenc}
\usepackage{mathptmx}
\usepackage{cite}
\usepackage{amsmath}
\usepackage{graphicx}
\usepackage{array}
\usepackage{booktabs}
\usepackage{multirow}
\usepackage{url}
\usepackage{hyperref}
\hypersetup{hidelinks}

\begin{document}

\title{SpendWise: AI-Powered Personal Finance Management - Expanded Implementation and Platform Analysis}

\author{
\IEEEauthorblockN{Lekhan HR}
\IEEEauthorblockA{\textit{Developer} \\
Bengaluru, India \\
lekhan.dev@email.com}
\and
\IEEEauthorblockN{Dr. Evelyn Reed}
\IEEEauthorblockA{\textit{HCI \& FinTech Research Lab} \\
Metropolitan Institute of Technology \\
e.reed@met.edu}
}

\maketitle

\begin{abstract}
Personal financial management remains difficult for many users because raw expense logs rarely translate into clear action. SpendWise addresses this problem by combining transaction tracking, budgeting, savings goals, analytics, notifications, and AI-guided financial interpretation in a unified software ecosystem. This expanded IEEE-style paper builds on the earlier SpendWise paper in the repository by incorporating the current implementation state of the codebase rather than describing only the earlier baseline mobile application. The present repository now includes a production-oriented web workspace, a maintained Expo React Native client, a parallel Flutter client, an auxiliary Expo extraction workspace, documentation assets, and continuous delivery workflows. In addition to the original core finance modules, the current codebase implements an investment console, browser-side import/export tooling, recurring bill detection, panel-based notifications, custom categories, PDF reporting, and other advanced utilities. This paper therefore serves both as an expanded technical description of the SpendWise platform and as a repository-grounded implementation report covering architecture, modules, data flow, automation, testing, limitations, and future expansion opportunities.
\end{abstract}

\begin{IEEEkeywords}
Personal finance, React, React Native, Expo, Flutter, Firebase, Groq, analytics, budgeting
\end{IEEEkeywords}

\section{Introduction}
SpendWise is designed to help users observe spending behavior, control category-level expenses, plan savings, and receive actionable feedback rather than only store raw financial entries. The repository reflects this goal through a layered implementation: user-facing clients, shared financial concepts, cloud-backed storage, AI inference, and deployment automation.

The current project is broader than a single mobile app. The \texttt{web/} directory contains the most feature-rich implementation, the \texttt{mobile/} directory contains the maintained Expo React Native client, the \texttt{flutter\_app/} directory contains a parallel Flutter implementation, and the repository also includes \texttt{expo-app/} as an extracted/mobile workspace plus \texttt{docs/} for published documentation. As a result, the project should be understood as a finance platform with multiple client surfaces rather than a single executable.

\section{Problem Statement and Objectives}
Most basic finance trackers stop at data entry: they let a user record income and expenses but do not actively help the user interpret spending behavior, anticipate risk, or convert data into short-term decisions. SpendWise addresses that gap through a combination of financial tracking, visual analytics, realtime cloud persistence, and AI-based interpretation.

From the implementation inspected in the repository, the project pursues the following concrete objectives:
\begin{enumerate}
\item provide a unified ledger for daily income and expense recording;
\item maintain per-category budget controls with utilization feedback and alerts;
\item support savings goals with deadlines, priority, and progress tracking;
\item generate interpretable financial analytics instead of raw lists only;
\item use AI to convert current data into health scores, summaries, and recommendations; and
\item deliver the experience across browser and mobile-oriented clients with synchronized cloud data.
\end{enumerate}

The expanded web workspace introduces an additional objective: to move from pure money tracking into financial planning by giving users entry-level investment guidance, market context, and roadmap-style next actions.

\section{Repository Composition and Scope}
Table~\ref{tab:repo} summarizes the major repository areas and their implemented roles.

\begin{table*}[t]
\caption{Major Repository Components}
\label{tab:repo}
\centering
\resizebox{\textwidth}{!}{%
\begin{tabular}{p{2.3cm}p{3.1cm}p{10.5cm}}
\toprule
\textbf{Path} & \textbf{Primary Stack} & \textbf{Implemented Role} \\
\midrule
\texttt{web/} & React 18, Vite 5, TypeScript, Firebase, Recharts, Capacitor & Main browser application with authentication, dashboard, transactions, budgets, reports, goals, settings, AI assistant, notification center, investment console, CSV/OFX export, PDF report generation, duplicate detection, recurring transaction tools, bills, achievements, and category management. \\
\texttt{mobile/} & Expo SDK 54, React Native 0.81, TypeScript, Firebase, react-native-chart-kit & Maintained mobile client with authentication, dashboard, transaction CRUD, budgets, goals, reports, AI insights, chatbot modal, settings, multi-currency support, and Expo notification scheduling. \\
\texttt{flutter\_} \newline \texttt{app/} & Flutter, Provider, Firebase, fl\_chart & Parallel native implementation of the core finance workflow: login, dashboard, transactions, budget, goals, reports, AI insights, settings, chatbot, and real-time Firebase synchronization. \\
\texttt{expo-} \newline \texttt{app/} & Expo, React Native, Firebase & Auxiliary extracted Expo workspace documenting a web-to-mobile conversion path and preserving an additional React Native implementation scaffold. \\
\texttt{docs/} & Static HTML/CSS/JS, PDFs & GitHub Pages documentation site, screenshots, architecture pages, and bundled academic/reporting assets. \\
\texttt{.github/} \newline \texttt{workflows/} & GitHub Actions & Continuous integration, Pages deployment, Vercel deployment, Android APK builds via EAS and Capacitor/Gradle, and test execution. \\
\bottomrule
\end{tabular}}
\end{table*}

The practical scope of the repository can therefore be divided into three layers:
\begin{itemize}
\item \textbf{Core finance layer:} transactions, budgets, goals, settings, categories, currency formatting, analytics, and AI summaries.
\item \textbf{Client delivery layer:} web, Expo React Native, Flutter, and documentation publishing.
\item \textbf{Operational layer:} authentication, persistence, notifications, testing, and deployment automation.
\end{itemize}

\section{System Architecture and Data Flow}
\subsection{Backend and Persistence}
The project uses Firebase Authentication and Firebase Realtime Database as the primary backend for the implemented finance workflow. Both the web and mobile React-based clients subscribe to user-scoped database paths using realtime listeners so that state changes propagate without manual refresh. The Flutter client follows the same pattern through stream subscriptions in its provider layer.

\begin{table}[t]
\caption{Primary Firebase Realtime Database Nodes}
\label{tab:firebase}
\centering
\begin{tabular}{p{2.2cm}p{4.9cm}}
\toprule
\textbf{Node} & \textbf{Purpose} \\
\midrule
\texttt{/users/\{uid\}/} \newline \texttt{transactions} & Stores income and expense entries including amount, category, method, date, description, and tags. \\
\texttt{/users/\{uid\}/} \newline \texttt{budgets} & Stores per-category budget limits, period, and notification preferences. \\
\texttt{/users/\{uid\}/} \newline \texttt{goals} & Stores savings goals with target amount, current amount, deadline, and priority. \\
\texttt{/users/\{uid\}/} \newline \texttt{settings} & Stores dark mode, currency, and in the web app also custom categories. \\
\bottomrule
\end{tabular}
\end{table}

\subsection{State Management}
The two React-based implementations center their logic around an \texttt{AppContext}. This context owns the finance entities, the active view, formatting helpers, CRUD actions, and persistence synchronization. The web client extends this pattern with a dedicated \texttt{NotificationContext} and \texttt{InvestmentContext} to isolate browser-only features such as investment tracking, market overview, and panel-based notifications.

\subsection{AI Insight Pipeline}
SpendWise integrates Groq's OpenAI-compatible chat completions endpoint to generate financial analysis. The implemented pipeline computes a monthly summary from transactions, derives category spending, budget consumption, and goal progress, and injects those aggregates into a prompt for the \texttt{llama-3.3-70b-versatile} model. The expected output is structured JSON containing a health score, insight cards, recommendations, and a short summary. If the API fails, deterministic fallback insights are generated locally from savings rate, budget use, and goal status.

\subsection{Notification Model}
The mobile client uses Expo Notifications for permission handling, budget alerts, daily reminders, and weekly summaries. The web client uses browser/native notifications through Capacitor plus a local notification center backed by \texttt{localStorage}. This split is important: mobile notifications are device-scheduled features, while the web notification center also acts as an in-app event feed for investment and planning workflows.

\subsection{Implemented Financial Metrics}
Several parts of the platform compute derived finance metrics rather than only display raw entries. The main formulas visible in the codebase can be summarized as follows:
\begin{equation}
\text{Savings Rate} = \frac{\text{Income} - \text{Expenses}}{\text{Income}} \times 100
\end{equation}
\begin{equation}
\text{Budget Utilization} = \frac{\text{Spent in Category}}{\text{Category Limit}} \times 100
\end{equation}
\begin{equation}
\text{Goal Progress} = \frac{\text{Current Saved Amount}}{\text{Target Amount}} \times 100
\end{equation}

These metrics drive the dashboard summary cards, reports, AI prompts, budget warnings, and goal pacing indicators. The web application goes further by forecasting month-end spend from current pace and by calculating monthly funding pressure for unfinished goals.

\subsection{Core Data Entities}
\begin{table*}[t]
\caption{Main Application Data Models}
\label{tab:entities}
\centering
\resizebox{\textwidth}{!}{%
\begin{tabular}{p{2.4cm}p{4.4cm}p{8.8cm}}
\toprule
\textbf{Entity} & \textbf{Main Fields} & \textbf{Implemented Role} \\
\midrule
Transaction & id, type, amount, category, paymentMethod, date, description, tags & Fundamental financial event model used by dashboards, reports, budgets, import/export, recurring detection, and AI summaries. \\
Budget & id, category, limit, period, startDate, notifications & Defines spending controls and supports percentage-based warning and over-limit logic. \\
Goal & id, name, targetAmount, currentAmount, deadline, priority & Tracks savings objectives and supports progress, pace, and deadline pressure calculations. \\
Category & id, name, icon, color, type & Binds visual identity and semantic grouping to transaction and budget workflows. \\
Currency & code, symbol, name, locale & Drives consistent formatting across clients and user preference persistence. \\
Notification & id, type, title, message, priority, timestamp, read & Used in the web notification center for panel-based event management and deduplication. \\
\bottomrule
\end{tabular}}
\end{table*}

\subsection{End-to-End User Workflow}
An end-to-end user session in SpendWise follows a consistent pattern across the principal implementations:
\begin{enumerate}
\item the user authenticates through Firebase;
\item the client subscribes to transactions, budgets, goals, and settings for that user;
\item the dashboard computes month-to-date metrics from local state hydrated by realtime listeners;
\item the user adds or edits transactions, budgets, or goals through forms and modal flows;
\item the updated arrays are written back to Firebase and reflected in all subscribed views;
\item analytics and AI modules consume those updated aggregates to produce charts, health scores, summaries, and recommendations; and
\item notification systems or reminder tools surface exceptions such as budget pressure, recurring bills, and planning milestones.
\end{enumerate}

\section{Implementation Methodology}
This expanded paper is based on repository-grounded inspection rather than on a standalone conceptual specification. The implementation methodology used for the report can be summarized in four steps:
\begin{enumerate}
\item \textbf{Manifest verification:} package manifests and configuration files were inspected to determine the active technology stack, framework versions, and deployment targets for each client surface.
\item \textbf{Source-level feature tracing:} component trees, context providers, service layers, and screen files were examined to identify which features are actually implemented and which are only documented.
\item \textbf{Cross-client comparison:} the web, Expo mobile, Flutter, and auxiliary Expo workspaces were compared to distinguish shared core functionality from client-specific extensions.
\item \textbf{Operational inspection:} GitHub Actions workflows, documentation assets, and export/testing utilities were reviewed to assess delivery readiness and engineering completeness.
\end{enumerate}

This approach matters because the repository contains both active implementations and older documentation snapshots. A repository-grounded method prevents outdated claims from being carried into the paper without verification.

\section{Implemented Functionalities}
\subsection{Cross-Platform Core Finance Modules}
The following finance modules are implemented across the main application surfaces:
\begin{itemize}
\item \textbf{Authentication:} Firebase email/password authentication is implemented across the major clients, while the web client also exposes Google sign-in and explicit email verification handling.
\item \textbf{Dashboard:} Users can inspect balance, income, expense, savings rate, recent activity, and summary cards immediately after login.
\item \textbf{Transaction management:} Users can create, edit, delete, search, and review financial entries with categories, descriptions, dates, payment methods, and tags.
\item \textbf{Budget control:} Users can create category budgets, monitor utilization, and view progress against spending limits.
\item \textbf{Savings goals:} Users can define goals with target amount, deadline, and priority and track progress over time.
\item \textbf{Analytics and reports:} Users can inspect visual breakdowns of spending trends over time through pie and bar-based reporting views.
\item \textbf{AI insights:} The project computes a financial health score and generates contextual insights and recommendations from current user data.
\item \textbf{Settings and preferences:} Currency selection, dark mode behavior, logout, and data reset capabilities are implemented.
\item \textbf{Realtime synchronization:} Firebase listeners keep state aligned between local UI state and the cloud data model.
\end{itemize}

\subsection{Detailed Behavior of Core Screens}
\subsubsection{Dashboard}
The web dashboard is not a static summary page. It computes month-to-date income, expense, balance, savings rate, average daily spend, days left in the month, budget usage, and goal funding status. It also renders a 14-day cashflow area chart, a spend-composition pie chart, budget pressure cards, goal progress cards, and a recent activity list. This makes the dashboard both a diagnostic overview and a navigation hub into more specialized modules.

The maintained mobile dashboard follows the same finance intent in a mobile-optimized form, while the Flutter client mirrors the same pattern through provider-driven state and native widgets. Together, these three implementations establish the dashboard as a shared conceptual core of the platform.

\subsubsection{Transaction Management}
Transaction handling is a major implemented subsystem rather than a single form. The web transaction screen supports free-text search, deferred filtering, type filtering, category filtering, date-window filtering, grouped daily views, edit/delete controls, category mix analysis, payment-method distribution, and largest-transaction highlights. Transactions also act as the source data for recurring bill detection, reporting, export, and AI context generation.

The mobile client includes transaction listing and analytics through its tabbed interface, while Flutter exposes similar CRUD logic through the provider and database service layers. This consistency confirms that transaction management is the system-of-record feature for the repository.

\subsubsection{Budget Management}
Budget management is implemented with per-category budget lines, visual progress bars, remaining-spend calculations, projected month-end spend, status labels such as on-track or over-limit, and identification of unbudgeted categories. On mobile, this same feature family is supported with charted or card-based summaries plus notification scheduling through Expo. In practical terms, the budget subsystem turns passive spending history into an active control mechanism.

\subsubsection{Goal Tracking}
Goal management is implemented with target amounts, current savings, deadline-aware pacing, priority labels, progress bars, and quick-add contribution actions in the web client. The code also computes how much funding is required per month to stay on schedule. This is an important distinction: the goal module is not only archival, but actively estimates whether a user is ahead of pace, behind pace, or overdue.

\subsubsection{Reports and Analytics}
The web reports screen extends beyond simple charts by offering range presets, cashflow comparison over the last six months, category contribution rankings, payment-method distribution, net-position cards, and CSV export. The mobile reports screen provides time-range switching with pie and bar charts. The Flutter implementation likewise includes a reports screen, reinforcing analytics as a first-class capability across the repository.

\subsubsection{Settings and Control Surface}
Settings are implemented as a control center rather than only a preference page. In the web client, the settings screen contains profile details, theme toggles, currency selection, custom category creation, data reset, import/export tools, PDF generation, and access to auxiliary tools such as savings calculator, receipt scanner, recurring transactions, category trends, achievements, and custom alerts. This screen effectively hosts the project's power-user workflow.

\subsection{AI Assistance Surfaces}
SpendWise implements AI in two related ways. First, it generates structured financial insights by computing a monthly summary and submitting it to Groq for JSON-formatted analysis. Second, it exposes conversational chatbot interfaces in the web, mobile, and Flutter-oriented clients. The web chatbot enriches its system prompt with the user's current financial context and allows the user to request pre-computed insights directly from the assistant interface.

\subsection{Detailed Web-Only Planning Extensions}
\subsubsection{Investment Workspace}
The investment module is implemented as a distinct product surface within the web client. Its current code includes a guided ``Start Here'' view, a market overview, stocks, portfolio, schemes, and roadmap tabs. The supporting context layer calculates risk bands from profile attributes such as horizon, experience, emergency fund depth, and objective; derives target allocation mixes; builds portfolio snapshots; and generates roadmap actions for rebalancing and goal de-risking. This module moves the application beyond bookkeeping into guided personal finance planning.

\subsubsection{Notification Center}
The web notification system maintains persisted notifications with read/unread state, priority, deduplication, dismissal, and retention limits. Unlike transient browser popups alone, this gives the project a recoverable in-app event feed for reminders and finance-related prompts.

\subsubsection{Import, Export, and Utility Tooling}
The browser workspace includes a substantial set of user utilities. CSV import supports preview and fuzzy duplicate detection. CSV and OFX export generate interoperable output files. Monthly PDF reporting is generated through a browser print flow. Bills and recurring transaction tools use stored data and heuristic grouping to detect repeating expenses. Achievements and category trends provide gamification and monitoring support. Receipt scanning exists as a prototype with simulated OCR and editable extracted values before conversion into transactions.

\subsection{Web-Specific Extended Modules}
The web client is presently the most advanced implementation in the repository. In addition to the shared finance workflow, it implements:
\begin{itemize}
\item \textbf{Enhanced reporting:} week/month/three-month/year filters, six-month cashflow comparison, category contribution breakdown, payment-method analysis, and CSV export from the reports screen.
\item \textbf{Investment console:} a dedicated investment workspace with guided onboarding, market overview, stock watchlists, portfolio tracking, schemes, roadmap actions, and profile-based allocation logic.
\item \textbf{Notification center:} read/unread state, grouped notifications, dedupe keys, and persistence with a capped in-app queue.
\item \textbf{Data import/export:} CSV import with preview, fuzzy duplicate detection, OFX export, and browser-generated PDF financial summaries.
\item \textbf{Category management:} creation and removal of custom income and expense categories directly from the settings flow.
\item \textbf{Utility modules:} savings calculator, category trends, achievements, recurring transactions, custom alerts, and bills/reminders.
\item \textbf{Receipt scanner prototype:} a browser upload flow with simulated OCR extraction. The code explicitly marks this as a mock implementation intended for later replacement by real OCR.
\end{itemize}

\subsection{Additions Beyond the Earlier IEEE Paper}
The earlier IEEE paper in the repository focused mainly on the baseline personal finance workflow and empirical evaluation storyline. The current repository has moved beyond that baseline in several important ways, summarized in Table~\ref{tab:additions}.

\begin{table*}[t]
\caption{Major Additions Beyond the Earlier SpendWise IEEE Paper}
\label{tab:additions}
\centering
\resizebox{\textwidth}{!}{%
\begin{tabular}{p{4.0cm}p{4.6cm}p{6.8cm}}
\toprule
\textbf{Area} & \textbf{Earlier Paper Focus} & \textbf{Expanded Current Implementation} \\
\midrule
Primary application surface & Cross-platform finance app baseline & Full browser workspace plus maintained Expo client, Flutter client, and auxiliary Expo extraction project \\
Analytics depth & Weekly/monthly/yearly charts and AI score & Six-month cashflow comparison, category contribution cards, payment-method analysis, export actions, and richer dashboard metrics \\
Planning features & Budgets and goals & Dedicated investment module with market overview, schemes, portfolio tracking, risk-band logic, and roadmap actions \\
Notification handling & Budget alerts and weekly summaries & Web notification center with dedupe, persistence, grouped state, and investment-related event handling \\
Data portability & Not emphasized & CSV import with duplicate detection, CSV export, OFX export, and browser-generated PDF reports \\
Utility surface & Core finance tabs & Bills, recurring transactions, achievements, category trends, custom categories, savings calculator, and receipt scanning prototype \\
Operational maturity & Research-paper framing & GitHub Actions CI, Vercel deployment, GitHub Pages publishing, EAS Android builds, and Capacitor APK workflow \\
\bottomrule
\end{tabular}}
\end{table*}

\subsection{Client Coverage}
Table~\ref{tab:coverage} summarizes how the current repository distributes features across client surfaces.

\begin{table*}[t]
\caption{Implemented Feature Coverage by Client}
\label{tab:coverage}
\centering
\resizebox{\textwidth}{!}{%
\begin{tabular}{p{4.2cm}ccccp{5.0cm}}
\toprule
\textbf{Feature Area} & \textbf{Web} & \textbf{Mobile} & \textbf{Flutter} & \textbf{Expo-App} & \textbf{Remarks} \\
\midrule
Firebase authentication & Yes & Yes & Yes & Yes & Web adds Google sign-in and explicit email verification flow. \\
Transactions, budgets, goals & Yes & Yes & Yes & Yes & Core CRUD workflow is implemented in all major app variants. \\
Dashboard and reporting & Yes & Yes & Yes & Partial & The web client provides the richest analytics view. \\
AI insights / chatbot & Yes & Yes & Yes & Partial & All main clients include AI-facing functionality; richness varies by surface. \\
Notifications & Yes & Yes & Partial & Partial & Mobile uses Expo scheduling; web mixes browser notifications and panel-only alerts. \\
Investment module & Yes & No & No & No & Implemented only in the web application. \\
CSV / OFX / PDF export & Yes & No & No & No & Implemented as browser utilities in settings and reports. \\
Recurring bills / reminders & Yes & No & No & No & Web uses \texttt{localStorage}-backed helpers and recurring detection. \\
Custom categories & Yes & No & No & No & Persisted as part of web settings. \\
\bottomrule
\end{tabular}}
\end{table*}

\section{Technology Stack, Automation, and Testing}
\subsection{Technology Stack}
\begin{table*}[t]
\caption{Implemented Technology Stack}
\label{tab:stack}
\centering
\resizebox{\textwidth}{!}{%
\begin{tabular}{p{2.6cm}p{3.9cm}p{9.8cm}}
\toprule
\textbf{Layer} & \textbf{Technology} & \textbf{Implemented Use} \\
\midrule
Web frontend & React 18.2, Vite 5.2, TypeScript, Recharts, Tailwind CSS, Capacitor 8 & Browser UI, component architecture, reports, native packaging path, and responsive interaction model. \\
Mobile frontend & Expo SDK 54, React Native 0.81.5, React 19.1, Expo Router, react-native-chart-kit, Expo Notifications & Cross-platform mobile navigation, charts, notifications, and device-oriented finance workflows. \\
Flutter frontend & Flutter, Provider, fl\_chart, shared\_preferences, flutter\_local\_notifications & Parallel native app implementation with provider-based state and chart rendering. \\
Backend services & Firebase Auth, Firebase Realtime Database & Authentication, user-scoped storage, and realtime synchronization. \\
AI services & Groq chat completions, \texttt{llama-3.3-70b-versatile} & Health score generation, insight cards, recommendations, and summaries. \\
DevOps and hosting & GitHub Actions, EAS Build, GitHub Pages, Vercel & Continuous integration, Android release workflows, documentation hosting, and web deployment. \\
\bottomrule
\end{tabular}}
\end{table*}

\subsection{Automation and Delivery}
The repository contains a practical automation story rather than only local scripts. The \texttt{ci.yml} workflow installs dependencies, builds the web app, and runs Vitest tests across a Node.js matrix. Separate workflows deploy the documentation site to GitHub Pages and the web application to Vercel. Android delivery is implemented in two ways: an Expo/EAS workflow for the maintained mobile client and a Capacitor/Gradle-based APK workflow for the web package path.

\subsection{Testing and Verification}
Automated tests are currently concentrated in the web workspace. The repository includes unit and UI tests for:
\begin{itemize}
\item fuzzy duplicate detection during CSV import,
\item the settings import-preview flow, and
\item recurring bill detection behavior.
\end{itemize}
This test set validates important utility features but does not yet amount to full end-to-end coverage for every client or every finance module.

\begin{table}[t]
\caption{Current Verification Coverage}
\label{tab:verification}
\centering
\begin{tabular}{p{2.2cm}p{4.8cm}}
\toprule
\textbf{Target} & \textbf{Implemented Verification} \\
\midrule
Web utilities & Vitest coverage for duplicate detection, settings import workflow, and bills UI behavior \\
Web delivery & CI workflow installs dependencies, builds the app, and runs tests on multiple Node versions \\
Documentation & GitHub Pages workflow publishes the static documentation bundle \\
Web deployment & Vercel workflow deploys preview/production builds using repository secrets \\
Android output & Expo/EAS and Capacitor/Gradle workflows automate APK generation paths \\
\bottomrule
\end{tabular}
\end{table}

\section{Current Status and Limitations}
The project is functionally substantial and already deployable, but several implementation realities should be stated clearly:
\begin{itemize}
\item \textbf{Feature asymmetry exists across clients.} The web application is ahead of the mobile and Flutter clients in breadth, especially for investment and advanced utility modules.
\item \textbf{Some utilities are intentionally local prototypes.} For example, the receipt scanner currently uses simulated OCR and browser-only preview logic rather than production OCR.
\item \textbf{Some documentation is stale relative to manifests.} Older markdown files cite earlier Expo/React Native versions than those currently declared in package manifests.
\item \textbf{Supabase is scaffolded but not central.} A Supabase client file and dependency exist in the web workspace, but the implemented persistence path remains Firebase-based.
\item \textbf{Testing depth is still uneven.} Web utilities have automated tests, while the mobile and Flutter clients rely more heavily on manual verification at the present stage.
\end{itemize}

\section{Conclusion}
SpendWise is not a toy tracker but a multi-surface finance platform with real implementation depth. The repository already demonstrates a complete finance loop: user authentication, realtime financial record keeping, budgeting, goal tracking, reporting, and AI-based guidance. The web application extends this with investment workflows and several practical data tools, while the mobile and Flutter clients establish strong cross-platform coverage for the core product. From a project-report perspective, the repository shows solid breadth in feature engineering, cloud integration, and delivery automation, with the clearest future work being feature consolidation across clients, replacement of simulated utility logic with production services, and broader automated test coverage.

\section*{Acknowledgment}
The authors acknowledge the open-source ecosystem around React, Expo, Flutter, Firebase, and related tooling, as well as the repository documentation that made cross-verification of the SpendWise platform feasible.

\begin{thebibliography}{00}
\bibitem{b1} Meta Platforms, ``React Native Documentation,'' official documentation.
\bibitem{b2} Expo, ``Expo Documentation,'' official documentation.
\bibitem{b3} Google, ``Firebase Authentication Documentation,'' official documentation.
\bibitem{b4} Google, ``Firebase Realtime Database Documentation,'' official documentation.
\bibitem{b5} Groq, ``OpenAI-Compatible API Documentation,'' official documentation.
\bibitem{b6} Flutter Team, ``Flutter Documentation,'' official documentation.
\bibitem{b7} Vercel, ``Vercel Documentation,'' official documentation.
\end{thebibliography}

\end{document}
